% ==============================================================================
% T0 Theory: Document 166
% Title: The Casimir Scale as Cosmic Link:
%        Casimir Effect, CMB and Hubble Constant in a xi-Power Ladder
% Author: Johann Pascher
% Date: March 2026
% References: Documents 025, 026, 091, 165
% ==============================================================================

\section*{Abstract}

This document shows that the Casimir effect, the cosmic microwave background (CMB)
and the Hubble constant $H_0$ are not independent phenomena on different scales,
but manifestations of the same $\xi$ geometry at different power levels.
The Casimir vacuum scale $L_\xi \approx 100\,\mu\text{m}$ (Doc.~091) and the
H$_0$ ratio $\hbar H_0 / m_e c^2 = (\pi/2)\cdot\xi^{10}$ (Doc.~165)
are linked by an exact algebraic bridge formula:

\begin{equation}
    \boxed{L_\xi \cdot \frac{H_0}{c}
    = \frac{\pi}{2}\cdot\left(\frac{15}{\pi^2}\right)^{\!1/4}
    \cdot \frac{m_e c^2}{k_B T_\text{CMB}} \cdot \xi^{41/4}}
    \label{eq:bridge_main}
\end{equation}

The exponent $41/4$ is identical to the one in Doc.~026 ($E_H = E_0 \cdot \xi^{41/4}$)
and is composed of the H$_0$ exponent $10$ and the Casimir exponent $1/4$.
Five seemingly unconnected scales -- from the sub-Planck length to the Hubble length --
form a gapless $\xi$-power ladder with nearly constant step spacing.

\clearpage

% ==============================================================================
\section{Background: Three Apparently Unconnected Phenomena}
\label{sec:background}
% ==============================================================================

\subsection{Casimir Effect and Vacuum Structure (Doc.~091)}

Doc.~091 shows that the energy density of the cosmic microwave background
is described by the T0 vacuum as:

\begin{equation}
    \rho_\text{CMB} = \frac{\xi \hbar c}{L_\xi^4}
    \label{eq:casimir_cmb}
\end{equation}

Numerically, the characteristic vacuum length scale comes out as
$L_\xi \approx 100\,\mu\text{m}$ -- precisely the scale at which the Casimir energy density
and the CMB energy density coincide:

\begin{equation}
    \rho_\text{Casimir}(d = L_\xi)
    = \frac{\pi^2 \hbar c}{240\,L_\xi^4}
    = \frac{\pi^2}{240\,\xi}\,\rho_\text{CMB}
    \label{eq:casimir_at_Lxi}
\end{equation}

Substituting \eqref{eq:casimir_cmb}, equation~\eqref{eq:casimir_at_Lxi} reproduces
the standard Casimir formula exactly -- $\xi$ cancels out.

\subsection{H\texorpdfstring{$_0$}{0} Ratio (Doc.~165)}

Doc.~165 shows:

\begin{equation}
    \frac{\hbar H_0}{m_e c^2} = \frac{\pi}{2}\,\xi^{10}
    \label{eq:h0_ratio}
\end{equation}

This yields $H_0 \approx 66.8$ km/s/Mpc -- within $0.9\%$ of the
Planck value $67.4$ km/s/Mpc, without free parameters.

\subsection{The Open Question}

The two results -- \eqref{eq:casimir_cmb} and \eqref{eq:h0_ratio} -- were
developed independently of each other and refer to entirely different
physical phenomena. Are they really independent? Or are they two
aspects of the same $\xi$ structure?

% ==============================================================================
\section{Derivation of the Bridge Formula}
\label{sec:derivation}
% ==============================================================================

\subsection{L\texorpdfstring{$_\xi$}{xi} from T\texorpdfstring{$_\text{CMB}$}{CMB}}

Combining the Casimir vacuum relation \eqref{eq:casimir_cmb} with the
Stefan--Boltzmann radiation formula
$\rho_\text{CMB} = (\pi^2/15)(k_B T_\text{CMB})^4/(\hbar c)^3$
gives $L_\xi$ as a function of $T_\text{CMB}$:

\begin{equation}
    L_\xi = \left[\frac{15\,\xi}{\pi^2}\right]^{1/4}
    \frac{\hbar c}{k_B T_\text{CMB}}
    \label{eq:Lxi_from_T}
\end{equation}

Numerically: $L_\xi = 100.24\,\mu\text{m}$
(deviation from the value obtained directly from the measured $\rho_\text{CMB}$: $0.11\%$).

\subsection{Multiplication of the Two Relations}

We multiply $L_\xi$ from \eqref{eq:Lxi_from_T} by $H_0/c$ from
\eqref{eq:h0_ratio}:

\begin{align}
    L_\xi \cdot \frac{H_0}{c}
    &= \left[\frac{15\xi}{\pi^2}\right]^{1/4}
       \frac{\hbar c}{k_B T_\text{CMB}}
       \cdot \frac{\pi}{2}\,\frac{m_e c^2\,\xi^{10}}{\hbar c}
    \nonumber\\[6pt]
    &= \frac{\pi}{2}\left[\frac{15}{\pi^2}\right]^{1/4}
       \frac{m_e c^2}{k_B T_\text{CMB}}
       \cdot \xi^{10 + 1/4}
    \nonumber\\[6pt]
    &= \frac{\pi}{2}\left[\frac{15}{\pi^2}\right]^{1/4}
       \frac{m_e c^2}{k_B T_\text{CMB}}
       \cdot \xi^{41/4}
    \label{eq:bridge_derivation}
\end{align}

This is formula~\eqref{eq:bridge_main}. The derivation is exact -- no approximation,
no free parameters.

\begin{keyresult}[Bridge Formula Casimir $\leftrightarrow$ H$_0$]
    \begin{equation}
        L_\xi \cdot \frac{H_0}{c}
        = \underbrace{\frac{\pi}{2}\left(\frac{15}{\pi^2}\right)^{1/4}}_{K_\text{geo} = 1.7441}
        \cdot
        \underbrace{\frac{m_e c^2}{k_B T_\text{CMB}}}_{= 2.000 \times 10^9}
        \cdot\; \xi^{41/4}
        \label{eq:bridge_boxed}
    \end{equation}
    Numerical verification: $L_\xi \cdot H_0/c = 7.241 \times 10^{-31}$
    (direct and algebraic evaluation identical to within $10^{-14}\%$).
\end{keyresult}

\subsection{The Exponent 41/4 as Sum of Two Known Exponents}

\begin{equation}
    \frac{41}{4} = \underbrace{10}_{\text{H}_0\text{ exponent (Doc.~165)}}
                 + \underbrace{\frac{1}{4}}_{\text{Casimir exponent}}
    \label{eq:exponent_decomp}
\end{equation}

The same exponent appears in Doc.~026 as $E_H = E_0 \cdot \xi^{41/4}$.
This means that the decomposition $41/4 = 10 + 1/4$ is not accidental but
reflects the physical structure -- $10$ describes the
coupling strength from the electron mass to the Hubble scale, $1/4$ describes
the Casimir scaling of the vacuum energy.

% ==============================================================================
\section{The Complete \texorpdfstring{$\xi$}{xi}-Power Ladder}
\label{sec:ladder}
% ==============================================================================

\subsection{Five Scales in One Ladder}

We compute all relevant scales as $\xi$ powers of the Planck length $\ell_P$:

\begin{table}[H]
    \centering
    \caption{The $\xi$-power ladder: five length scales and their $\xi$ exponents}
    \label{tab:scale_ladder}
    \adjustbox{max width=\textwidth}{%
\begin{tabular}{lccc}
        \toprule
        \textbf{Scale} & \textbf{Physical meaning}
            & \textbf{Length [m]} & \textbf{$n$: scale $\sim \xi^n \cdot \ell_P$} \\
        \midrule
        $L_0 = \xi \cdot \ell_P$ & Sub-Planck, T0 lattice & $2.16 \times 10^{-39}$ & $n = +1.0$ \\
        $\ell_P$                  & Planck length           & $1.62 \times 10^{-35}$ & $n = \phantom{+}0.0$ \\
        $\bar\lambda_e$           & Compton wavelength      & $3.86 \times 10^{-13}$ & $n = -5.77$ \\
        $L_\xi$                   & Casimir vacuum scale    & $1.00 \times 10^{-4}$  & $n = -7.95$ \\
        $c/H_0$                   & Hubble length           & $1.38 \times 10^{26}$  & $n = -15.72$ \\
        \bottomrule
    \end{tabular}%
}
\end{table}

\subsection{Uniform Step Structure}

\begin{table}[H]
    \centering
    \caption{Step spacings in the $\xi$-power ladder}
    \label{tab:steps}
    \adjustbox{max width=\textwidth}{%
\begin{tabular}{lcc}
        \toprule
        \textbf{Transition} & \textbf{$\Delta n$} & \textbf{Factor} \\
        \midrule
        $L_0 \to \ell_P$        & $-1.0$ & $1/\xi$            \\
        $\ell_P \to \bar\lambda_e$ & $-5.77$ & $\xi^{-5.77}$ \\
        $\bar\lambda_e \to L_\xi$ & $-2.18$ & $\xi^{-2.18}$ \\
        $L_\xi \to c/H_0$        & $-7.78$ & $\xi^{-7.78}$ \\
        \midrule
        $\ell_P \to c/H_0$ (total) & $-15.72$ & $\xi^{-15.72}$ \\
        \bottomrule
    \end{tabular}%
}
\end{table}

\subsection{The Geometric Mean}

A remarkable finding: $L_\xi$ lies almost exactly at the
\emph{geometric mean} of $\ell_P$ and $c/H_0$:

\begin{equation}
    \sqrt{\ell_P \cdot \frac{c}{H_0}} = 47.3\,\mu\text{m},
    \qquad
    L_\xi = 100.2\,\mu\text{m},
    \qquad
    \frac{L_\xi}{\sqrt{\ell_P \cdot c/H_0}} = 2.12
    \label{eq:geom_mean}
\end{equation}

\begin{insightBox}[The Casimir Scale as Cosmic Link]
    On the $\xi$-power ladder, the Casimir vacuum scale $L_\xi \approx 100\,\mu\text{m}$
    lies precisely between the Planck length and the Hubble length.
    It is neither arbitrary nor accidental, but geometrically determined
    by the same parameter $\xi$ that generates all the other scales.
\end{insightBox}

% ==============================================================================
\section{Physical Interpretation}
\label{sec:interpretation}
% ==============================================================================

\subsection{Prefactors and Their Meaning}

The bridge formula \eqref{eq:bridge_boxed} contains two dimensionless factors:

\begin{description}
    \item[$K_\text{geo} = (\pi/2)(15/\pi^2)^{1/4} = 1.744$]
        A purely geometric factor arising from 3D Stefan--Boltzmann radiation
        (dimension of phase space) and the torus topology of the T0 time field
        (factor $\pi/2$, cf.\ Doc.~159).

    \item[$m_e c^2 / (k_B T_\text{CMB}) = 2.00 \times 10^9$]
        The ratio of the electron rest energy to the CMB thermal energy.
        That this is a smooth number $\approx 2 \times 10^9$ shows
        the deep connection between the electroweak scale and the
        cosmological temperature.
\end{description}

\subsection{Why the Exponent 41/4 Is Fundamental}

The decomposition $41/4 = 10 + 1/4$ has the following meaning:

\begin{itemize}
    \item The exponent $10$ in $\hbar H_0/(m_e c^2) = (\pi/2)\xi^{10}$
    describes how far the cosmological energy scale lies from the
    particle mass scale -- measured in units of $\xi$.

    \item The exponent $1/4$ in $L_\xi \propto \xi^{1/4}$ describes
    how the Casimir scale grows away from the Planck scale -- as a
    fourth root, because the vacuum energy scales as $1/L^4$.

    \item The sum $41/4$ appears in Doc.~026 as the exponent in
    $E_H = E_0 \cdot \xi^{41/4}$. The present derivation explains
    \emph{why}: it is the sum of the two natural exponents
    of the Casimir relation and the H$_0$ relation.
\end{itemize}

\begin{keyresult}[Explanation of the Exponent 41/4]
    The exponent $41/4$ found empirically in Doc.~026 in
    $E_H = E_0 \cdot \xi^{41/4}$ is not arbitrary. It arises
    as the sum:

    \begin{equation}
        \frac{41}{4}
        = \underbrace{10}_{\hbar H_0/(m_e c^2) \,\sim\, \xi^{10}}
        + \underbrace{\frac{1}{4}}_{L_\xi \,\sim\, \xi^{1/4} \cdot \ell_P}
        \label{eq:41_4_explained}
    \end{equation}

    It encodes the coupling between the vacuum structure (Casimir, $1/4$)
    and the cosmological scale (H$_0$, $10$).
\end{keyresult}

% ==============================================================================
\section{Experimental Predictions}
\label{sec:predictions}
% ==============================================================================

\subsection{Casimir Precision Measurements}

The T0 theory predicts that at a plate separation $d = L_\xi \approx 100\,\mu\text{m}$
the Casimir energy density coincides with the CMB energy density. Since $\rho_\text{CMB}$
is known very precisely, this yields a precise prediction for $L_\xi$:

\begin{equation}
    L_\xi^\text{T0}
    = \left(\frac{\xi\hbar c}{\rho_\text{CMB}}\right)^{1/4}
    = 100.24\,\mu\text{m}
    \label{eq:Lxi_prediction}
\end{equation}

This is a testable prediction: at $d = L_\xi$, Casimir measurements should
show a characteristic signature linked to the CMB energy
density.

\subsection{H\texorpdfstring{$_0$}{0} Measurement as Indirect Casimir Test}

The bridge formula \eqref{eq:bridge_main} makes a consistent prediction:
if $L_\xi$ is determined from Casimir measurements and $T_\text{CMB}$ is known,
$H_0$ follows directly:

\begin{equation}
    H_0 = \frac{\pi}{2}\left[\frac{15}{\pi^2}\right]^{1/4}
          \frac{m_e c^2}{k_B T_\text{CMB}}
          \cdot \frac{c}{L_\xi} \cdot \xi^{41/4}
    \label{eq:H0_from_Casimir}
\end{equation}

A Casimir experiment thus indirectly determines the Hubble constant --
via the same geometric parameter $\xi$ that connects the two.

% ==============================================================================
\section{Complete Numerical Chain}
\label{sec:numerical}
% ==============================================================================

\begin{table}[H]
    \centering
    \caption{Numerical verification of the complete chain}
    \label{tab:numerical}
    \adjustbox{max width=\textwidth}{%
\begin{tabular}{lll}
        \toprule
        \textbf{Quantity} & \textbf{T0 formula} & \textbf{Numerical value} \\
        \midrule
        $\xi$                      & $4/30000$                       & $1.333 \times 10^{-4}$ \\
        $L_\xi$                    & $[15\xi/\pi^2]^{1/4}\,\hbar c/(k_B T)$
                                                                      & $100.24\,\mu\text{m}$ \\
        $\rho_\text{CMB}$          & $\xi\hbar c/L_\xi^4$            & $4.170 \times 10^{-14}\,\text{J/m}^3$ ✓ \\
        $T_\text{CMB}$             & Stefan--Boltzmann               & $2.725\,\text{K}$ ✓ \\
        $\hbar H_0/(m_e c^2)$      & $(\pi/2)\,\xi^{10}$             & $1.896 \times 10^{-38}$ \\
        $H_0$                      & from the row above              & $66.82\,\text{km/s/Mpc}$ \\
        $L_\xi \cdot H_0/c$        & $K_\text{geo} \cdot (m_e c^2/k_BT) \cdot \xi^{41/4}$
                                                                      & $7.241 \times 10^{-31}$ ✓ \\
        \bottomrule
    \end{tabular}%
}
\end{table}

\begin{foundation}[Three Independent Inputs -- One Consistent System]
    From only three measured quantities ($\xi$, $T_\text{CMB}$, $m_e$), both the
    Casimir scale $L_\xi$ and the Hubble constant $H_0$ follow without any further
    free parameters. The bridge formula \eqref{eq:bridge_main}
    is satisfied exactly -- no approximation, no fit parameter.
\end{foundation}

% ==============================================================================
\section{Summary}
\label{sec:summary}
% ==============================================================================

\begin{foundation}[Results of This Document]
    \begin{enumerate}
        \item The Casimir scale $L_\xi \approx 100\,\mu\text{m}$ and the
              Hubble constant $H_0$ are connected by the exact formula
              $L_\xi \cdot H_0/c = K_\text{geo} \cdot (m_e c^2/k_B T_\text{CMB}) \cdot \xi^{41/4}$.
        \item The exponent $41/4 = 10 + 1/4$ explains the value found
              empirically in Doc.~026: $10$ from the H$_0$ scale,
              $1/4$ from the Casimir scaling.
        \item Five scales -- sub-Planck, Planck, Compton, Casimir, Hubble --
              form a $\xi$-power ladder with exponent steps
              $\Delta n \approx \{1;\,5.77;\,2.18;\,7.78\}$.
        \item The Casimir scale $L_\xi$ lies near the geometric mean
              of $\ell_P$ and $c/H_0$, with a factor of $2.12$.
        \item A Casimir precision experiment at $d \approx 100\,\mu\text{m}$
              indirectly tests the T0 prediction for $H_0$.
    \end{enumerate}
\end{foundation}

\vspace{1.5em}
\noindent\textbf{Related Documents:}
Doc.~025/026 (T0 cosmology, $E_H = E_0\xi^{41/4}$),
Doc.~091 (Casimir--CMB connection),
Doc.~159 (torus topology),
Doc.~165 ($H_0$ ratio $\xi^{10}$).
